% ===================================================================
\section{Introduction}
\label{ch:intro}
% ===================================================================

\subsection{Background and Motivation}

The Foundational Reasoning and Feedback Protocol (FRFP) is a framework for
structuring Human-AI collaboration through a typed pipeline architecture that
enforces separation between explicit (AI-accessible) and tacit
(human-exclusive) cognitive objects. The framework was published in two companion
papers: a conceptual/theory paper establishing the mathematical foundations, and a
technical specification paper defining the protocol semantics.

Informal mathematical proofs in papers --- even rigorous ones --- routinely omit
steps deemed ``obvious.'' These omissions are harmless in most cases, but in a
framework whose explicit goal is to provide formal governance guarantees for
AI-assisted decision-making, any gap in the mathematical foundation is a direct
threat to the stated claims.

Formal verification addresses this by requiring every logical inference step to be
checked by a proof assistant. When a proof assistant accepts a proof, it provides
the same level of certainty as a line-by-line expert review with no gaps allowed.

This paper reports a complete machine-checked formalization of FRFP using Lean~4,
the most active modern interactive proof assistant. The formalization was conducted
as an independent companion effort: no changes were made to the published papers,
but the Lean proof obligations revealed refinements, corrections, and new results
that are documented here.

% -------------------------------------------------------------------
\subsection{What Formal Verification Provides}
% -------------------------------------------------------------------

\textbf{Lean~4 guarantees when it accepts a proof:}

\begin{itemize}
  \item \textbf{Logical soundness}: Every inference step follows from axioms and
        previously proven results.
  \item \textbf{Completeness}: No gaps or hand-waving --- every ``trivial'' step
        must be justified.
  \item \textbf{Type correctness}: All variables carry consistent types through
        the entire proof.
  \item \textbf{Termination}: Proofs contain no circular reasoning.
\end{itemize}

\textbf{What Lean does not guarantee:}

\begin{itemize}
  \item That the formalized problem matches the intended informal problem.
  \item That axioms correspond to physical or computational reality.
  \item That the framework design is optimal or unique.
\end{itemize}

These caveats are standard in formal verification and define the scope of what
machine-checking can provide.

% -------------------------------------------------------------------
\subsection{Scope of This Companion Paper}
% -------------------------------------------------------------------

This companion covers:

\begin{enumerate}
  \item \textbf{21 core modules} comprising the full FRFP formal library
        (\texttt{Frfp/Core/}).
  \item \textbf{3 minimal-basis modules} (\texttt{Frfp/Minimal/}) implementing
        the reduced axiom basis.
  \item \textbf{1 case study module} (\texttt{Example/Sepsis.lean}) demonstrating
        an end-to-end executable pipeline.
  \item \textbf{1 executable report} (\texttt{FRFPReport.lean}) summarizing all
        key theorems.
\end{enumerate}

The verification is complete under two explicit caveats:
\begin{enumerate}
  \item The seven irreducible FRFP base axioms (HEG, HEC, CP, IL, AR, NTER, CSC)
        are assumed as theory premises.
  \item Every other axiom cites a published, research-grade reference.
\end{enumerate}

Under these caveats: \textbf{0 sorry, 0 uncited axioms, 0 build errors.}

% -------------------------------------------------------------------
\subsection{Primary Contributions}
% -------------------------------------------------------------------

\begin{contribution}[Complete machine-checked formalization]
269 theorems across 21 modules, verified by Lean~4.29.0 with Mathlib v4.29.0.
This is the first fully machine-checked formalization of a human-AI governance
framework at this scale.
\end{contribution}

\begin{contribution}[Axiom minimization]
The original FRFP basis of 11 axioms is reducible to 7 independent primitives.
Four axioms --- ETS, RB, AE, and MD --- are derivable as theorems from the
remaining core.
\end{contribution}

\begin{contribution}[New theorems]
Three genuinely novel results emerged from the proof obligations imposed by Lean:
\begin{itemize}
  \item A canonical earliest-unresolved-witness theorem for governance traces.
  \item A governance-safe trace-extension theorem.
  \item A protocol phase-decomposition (normal-form) theorem.
\end{itemize}
None of these were conjectured or stated in the published papers.
\end{contribution}

\begin{contribution}[Vibecoding methodology documentation]
The formalization was conducted using an AI-assisted development workflow. This
paper documents the workflow, its success and failure rates, and lessons for
others attempting similar formalization projects.
\end{contribution}

% -------------------------------------------------------------------
\subsection{Relation to the Published FRFP Papers}
% -------------------------------------------------------------------

\begin{center}
\begin{tabular}{ll}
\toprule
Paper & Role \\
\midrule
FRFP conceptual/theory & Defines framework; cited as the informal specification \\
FRFP technical specification & Defines protocol details; cited as the engineering reference \\
This paper & Machine-checked formalization; reports new structural results \\
\bottomrule
\end{tabular}
\end{center}

% -------------------------------------------------------------------
\subsection{Organization of This Paper}
% -------------------------------------------------------------------

\begin{center}
\begin{tabular}{ll}
\toprule
Chapter & Content \\
\midrule
2 & AI-assisted Lean development methodology (vibecoding) \\
3 & Axiom architecture, audit protocol, and citation inventory \\
4 & Core verification results by module \\
5 & Reduced axiom basis: 11 $\to$ 7 primitives \\
6 & New theorems from formalization \\
7 & Case studies: Sepsis example, Grothendieck construction, $\mathrm{runTo}\omega$ bijection \\
8 & Reproducibility package: build instructions, theorem inventory, module map \\
\bottomrule
\end{tabular}
\end{center}

% -------------------------------------------------------------------
\subsection{Availability}
% -------------------------------------------------------------------

The full Lean source code, build scripts, and documentation are available in the
\texttt{FRFP\_Math\_Verification/} repository. Chapter~\ref{ch:repro} provides a
complete guide to reproducing all results from a clean environment.
